\section{Equivalence Properties}
\label{sec:equivalence}

All three slice constructions produce identical results for every valid
input.

\begin{itemize}
  \item \href{https://github.com/thiagomata/prime-numbers/blob/list-article-v1.0.0/src/main/scala/v1/chapter3/list/properties/SliceEquivalenceLemmas.scala}{Slice equivalence}:
    tail-recursive, head-recursive, and index-range slices are identical
    for all valid inputs
\end{itemize}

\subsection{Slice Equivalence Lemma}

All three slice implementations --- tail-recursive, head-recursive, and
index-range --- produce the same result for any valid input.

\begin{gather*}
\forall\, L \in \mathbb{L},\ \forall\, i, j \in \mathbb{N},\ i \leq j < |L| \\
\operatorname{slice}(L, i, j) = \operatorname{headRecursiveSlice}(L, i, j) = \operatorname{indexRangeValues}(L, i, j)
\end{gather*}

\textbf{Proof by induction on $j - i$}\par\nopagebreak

\textbf{Base case}: $i = j$ --- all three produce $L_i :: L_e$.\par\nopagebreak

\textbf{Inductive step}: For $i < j$, each function decomposes into a
head element plus a recursive call on $(i+1, j)$ or $(i, j-1)$. By the
inductive hypothesis, the recursive calls produce equal sublists, and the
head element is the same, so the results are equal.

{\raggedright
This property is verified in
\href{https://github.com/thiagomata/prime-numbers/blob/list-article-v1.0.0/src/main/scala/v1/chapter3/list/properties/SliceEquivalenceLemmas.scala}{\texttt{SliceEquivalenceLemmas::\allowbreak tailHeadAndIndexRangeSlicesAreEqual}}.
The full Scala verification code is in Appendix~A.24.
\par}
